\documentclass[11pt]{article}
\usepackage[a4paper,margin=1in]{geometry}
\usepackage{amsmath,amssymb}
\usepackage{booktabs}
\usepackage{hyperref}
\usepackage{enumitem}
\usepackage{microtype}
\usepackage[numbers]{natbib}

\title{Conflux: Principal-Context Security, Evaluation and Verification for LLM Agents\\\large Work-in-progress fourth-year paper}
\author{Raya Buckley}
\date{31 July 2026}

\begin{document}
\maketitle

\begin{abstract}
Large-language-model agents combine information controlled by different
principals while performing externally visible actions. Previous work
introduced Influence Tracking with Extrapolated Security (ITES) and the
System-Level Evaluator for Defences (SLED). Conflux develops those ideas into a
provider- and benchmark-independent research framework. The current repository
has one canonical, fail-closed security kernel with immutable provenance,
independent read, authorisation, visibility and consent decisions, versioned
decision evidence, bounded explicit-state verification, authenticated dynamic
planning, and an offline result-producing CLI. Architectural repair and the
credential-free vertical slice are tested, but real-model, external-benchmark
and solver-binary claims remain gated on reproducible result artefacts. The
claim under test is containment: an action is allowed only when every
influencing principal's own permissions allow it, and an input fetched from
elsewhere cannot acquire authority it was never granted. Three independent
bounded analyses agree on all eleven planned variants, ten of which are refuted,
over ten distinct bodies and nine distinct refuted bodies, since two of the
variants differ in name only; the same three analyses agree on all 512 combinations of nine independent defects,
511 of them refuted;
the archived prototype reproduces its published trace count under a stated
enumeration change; and property reduction keeps its documented scaling
advantage on the fixtures exercised here. This paper distinguishes archived
evidence, current implementation evidence and planned experiments; numerical
placeholders are filled only by generated artefacts.
\end{abstract}

\section{Introduction}
LLM agents process requests, retrieved documents and tool output controlled by
principals with different authority. An external author should not acquire the
initiating Principal's permissions merely by influencing a model context.
Model-level prompt-injection defences improve instruction following, while
system-level approaches constrain effects externally
\citep{camel,designpatterns,struq,spotlighting}.
Model-level defences are genuine empirical security controls that reduce the
probability of inappropriate decisions; Conflux provides a complementary
worst-case authority-confinement layer whose property holds under arbitrary
model proposals.

For Principal Context $PC$ and action $a$, the restricted ITES rule is
\[
  \operatorname{Allow}(PC,a) \Longleftrightarrow
  PC \neq \varnothing \land \forall p \in PC,\;P(p,a).
\]
The explicit non-empty condition prevents vacuous authority. Provenance grows
monotonically, while organisational authority is supplied by an independent
policy decision point. An externally fetched object retains the authenticated
provenance of its actual source(s); it does not inherit the requesting user's
authority merely because the request was made on the user's behalf. For
ordinary derived objects, $PC(\text{output}) \supseteq PC(\text{execution
inputs})$; for scheduled executions, $PC(\text{scheduled})$ includes
$PC(\text{scheduling context})$. Current Conflux execution does not reduce
Principal Context; any future exception requires a separately accepted
semantics and proof obligation.

This paper contributes four things. First, a fail-closed security kernel with
one tested semantics: an action is allowed only when every influencing
principal's own permissions allow it, provenance grows monotonically, and no
proposal can assert or narrow its own context. Second, an argument-aware
disclosure and attribution design in which visibility and consent are separate
decisions, so a release is a recorded decision rather than a side effect.
Third, bounded evidence over that kernel: an explicit-state evaluator, a
property-reduced verification path with two independent backends, and a
reproduction of the archived prototype under a stated enumeration change.
Fourth, a reproduction discipline: every number this paper states is generated
from a retained bundle whose protocol and per-cell results are checked in the
repository, and the rows without such evidence say so.

\section{Previous-year contribution and archived evidence}
The previous prototype accumulated authorship provenance and intersected
permissions. Its bounded SLED results apply to that implementation and model.
The archived source and PDF of the previous-year artefact remain in
\texttt{research/publications/paper/}. \texttt{ARCHIVED.md} records the commit
that imported them and states that the submission provenance is unverified, and
\texttt{ARCHIVE\_MANIFEST.json} holds a SHA-256 hash for every pre-existing
file, which repository validation checks, so the archive cannot be edited
silently. Those files describe the previous prototype and its bounded results;
they are not evidence for the refactored kernel. Current experiments use
separately named legacy-reproduction and canonical suites.
This paper is a separate artefact and does not reuse those numbers: its
results are counted only from the retained bundles under
\texttt{research/output/runs/}, each of which carries the protocol, per-cell
result and trace records that the tables check. Where this paper refers to the
earlier results it cites the previous manuscript, and the results sections
below derive only from the current bundles; only the bounded SLED statements
above name the archived prototype, and they are labelled as applying to it.

\section{Related work and positioning}
ITES is structurally analogous to Biba low-water-mark contamination, enriched
with authenticated principal provenance and authority derived from the
organisation's existing access-control system.

Classical information-flow work supplies the vocabulary Conflux uses. Denning's
lattice model orders the classifications that a flow policy
compares~\citep{denning1976}; Biba's integrity considerations add the dual
obligation that low-integrity input must not reach a high-integrity
sink~\citep{biba1977}, and LOMAC applies low-water-mark integrity protection to
commodity operating systems~\citep{lomac}. The decentralized label model lets
each principal own the policy attached to the data it
creates~\citep{myersliskov}, and the language-based survey collects the
techniques that enforce such policies statically~\citep{sabelfeld2003}. Because
real systems must disclose something, the declassification literature names the
dimensions along which a policy may be relaxed and the principles that keep the
relaxation safe~\citep{declassificationsurvey}; robust declassification demands
that an adversary cannot exploit the relaxation
itself~\citep{robustdeclassification}; nonmalleable information flow control
closes the composition gap that appears when several relaxed components
meet~\citep{nonmalleableifc}; and attacker control and impact make the
comparison quantitative by asking what an adversary can still
influence~\citep{attackercontrol}. Conflux keeps that vocabulary for its
decision records, but its principals are the parties in an agent session, its
labels are provenance and consent records, and a disclosure is an explicit
recorded decision rather than a byte-level flow.

Classical IFC systems
including HiStar~\citep{histar}, Flume~\citep{flume}, and
Asbestos~\citep{asbestos} demonstrate that trusted reference-monitor
boundaries, dynamic contamination, and multi-principal isolation are
well-established systems problems. Clark-Wilson~\citep{clarkwilson} provides
a model of integrity through certified transformations and separation of
duties, framing the question of under what explicitly modelled operation
conservative influence may be reduced. Conflux does not currently
implement an endorsement or trusted-transformation mechanism.
Wu, Cecchetti, and Xiao~\citep{fsecure} prevent untrusted inputs from
reaching privileged planning through a structured IFC pipeline; Conflux
instead treats all model-visible information as potentially influential and
derives authority from the intersection of all influencers' ACS permissions.

CaMeL isolates control and data flows with a programmable policy interface
(Python functions over tool name and arguments, with capability/source
tracking and a STRICT dependency mode), while privilege-oriented systems such as
Progent mediate tool calls with symbolic policy \citep{camel,progent}. PACT
tracks argument-level provenance and FORGE enforces history-sensitive Datalog
policies \citep{pact,forge}. AgentDojo supplies realistic prompt-injection tasks
and metrics \citep{agentdojo}. Conflux therefore does not claim priority for
system-level mediation, provenance or privilege control. Its narrower research
focus is collective Principal Context derived from authenticated provenance,
interpreted through organisational policy, and evaluated with a shared
operational and verification transition model. A defence-native property that
does not imply PE is a non-implication between objectives, not a defect in the
compared system.

\section{Canonical architecture}
Immutable domain values represent Principals, Principal Contexts, permissions,
resources, artefacts, provenance, sessions, actions and environment snapshots.
Policy ports decide read access, authorisation, visibility and consent
independently. The ITES kernel is pure: model proposals cannot assert or narrow
their Principal Context, alternative branches share an immutable parent, and
effect execution requires the exact decision certificate.

The code keeps one dependency direction: immutable domain values, policy and
provider ports that declare protocols without choosing implementations, an
application layer that composes decisions, and adapters that translate external
systems without redefining security decisions. The kernel owns the only
mediation boundary, and planning routes every effect back through it at action
time rather than at plan time. Verification is separate from execution: a
callback-free intermediate representation, a reference interpreter, and optional
solver backends that never replace the operational exploration. Open-ended plans
use authenticated operation catalogues, typed bindings, explicit loops and
continuations, and visible resource bounds, so a plan cannot name an operation
the catalogue does not authenticate.

\subsection{Proposal semantics}
A versioned proposal batch declares either \texttt{ALTERNATIVES} or
\texttt{ORDERED\_PLAN}. Alternatives are independently explored from one
parent. Ordered steps propagate state but are re-authorised immediately before
execution; no batch grants future authority.

\subsection{Evidence and failure semantics}
Reports distinguish proposed, authorised, blocked, executed, provider-failed
and incomplete actions. Policy errors, malformed proposals, missing provenance
and exhausted bounds are explicit fail-closed outcomes. Securely blocked
attacks are diagnostics, not executed invariant failures.

\section{Implementation status at this snapshot}

\begin{tabular}{p{0.24\linewidth}p{0.32\linewidth}p{0.35\linewidth}}
\toprule
Area & Implemented boundary & Evidence limitation \\
\midrule
Security kernel & Immutable domain values, independent policy ports, one ITES
kernel, isolated alternatives, ordered plans, and certificate-bound execution
& Assumes complete mediation, authenticated provenance, and faithful policy
adapters \\
Argument and disclosure & Trusted action-argument roles, pointwise selector
policy, audience-specific disclosure, deterministic redaction, and structured
attribution & Retained native mutants exercise finite monitor models, not all
deployment observers \\
Runtime and CLI & Strict scenarios, deterministic providers and traces,
reporting, chat boundary, planning demo, SLED and formal commands & Live
effects and models remain optional; filesystem writes default to dry-run \\
Native verification & Breadth-first SLED with canonical state, explicit bounds,
deduplication, retained bound states, and shortest counterexamples & Verdicts
apply only to the finite operational model and declared properties \\
Verification reduction & Property-scoped cone-of-influence closure, stable rule
identity, reference comparison, and counterexample lifting & Retained evidence
covers checked finite fixtures; partial-order and symmetry reduction remain
future work \\
Dynamic planning & Authenticated catalogues, typed bindings, immutable patches,
bounded loops and continuations, subplans, mediated code capabilities, outcome
validation, and authority-minimising selection & Comparative real-model utility
evidence is not yet retained \\
Formal subset & Serializable IR, reference interpreter, runtime differential
tests, optional Z3 BMC, and nuXmv Boolean subset & Solver binaries and unbounded
runtime equivalence are not established \\
AgentDojo & Exact package 0.1.35 / benchmark v1.2.2 translation with preserved
IDs, native metrics, failure categories, and raw fixture & The fixture proves
translation only; a live no-defence-versus-ITES comparison is pending \\
Policy and delegation & Optional pinned Cedar adapter plus exact, expiring,
revocable, one-use delegation records and disabled runtime consumption & Cedar
preflight does not establish live parity; delegation transfers no runtime
authority \\
\bottomrule
\end{tabular}

This table states repository capabilities rather than empirical efficacy. The
claim ledger, task registry, schemas, tests, and retained artefacts are the
sources of current evidence. Archived reports and the previous paper record
project history but do not define this snapshot.

\section{Evaluation methodology}
The deterministic semantic corpus covers empty and mixed contexts,
author/reader inversions, derivation, nested accumulation, branch isolation,
ordered plans, consent, visibility, unsupported delegation, policy changes and
provider failure. Executable defective variants must yield minimal
counterexamples.

Native SLED uses versioned legacy-reproduction and corrected canonical suites.
SLED-MC performs breadth-first exploration with canonical state keys,
deduplication, predecessor edges and explicit \texttt{SAFE},
\texttt{BOUNDED\_SAFE}, \texttt{UNSAFE} and \texttt{UNKNOWN} verdicts.

Real-model manifests record model revision, endpoint configuration, sampling,
structured-output mode, raw-response hashes, token use and latency. AgentDojo
integration pins an upstream revision, preserves upstream IDs, and reports
native and Conflux metrics without hiding added provenance assumptions.

AgentDojo breadth uses the frozen stratified selection: four pinned suites
(banking, slack, travel, workspace), each executed as one protocol whose suite
record names the upstream tag, the selected task ids and the injection task
ids. Every cell runs the same model configuration with seed and temperature
fixed, and each retained bundle keeps its protocol record, its per-cell result
JSON and its raw upstream traces. Repeatability is measured, not assumed: a
retained protocol is re-executed in its own directory and compared cell by
cell on status, native utility, native security, model-call count and failure
list, with the raw traces compared after the two timing fields are set aside,
and the protocol fingerprint must reproduce exactly. A repeated run is restored
byte-for-byte afterwards, so the retained bundle stays the artefact the tables
count. Failures are classified from the message that caused them, and each
cell keeps a bounded summary of that message, so an incomplete cell is
diagnosable from the bundle itself. At the retained floor-model scale utility
is zero, so the security counts measure the harness rather than a defence.

\section{Formal verification and planning}
SLED-V accepts a finite serialisable transition IR rather than arbitrary
Python. Initial properties cover unauthorised effects and reads, provenance
monotonicity, absence of implicit delegation and bounded resource use. Bounded
SMT results retain solver artefacts; unbounded claims require an
IC3/PDR-capable backend and runtime-to-IR conformance.

Property-scoped cone-of-influence reduction closes transitively over guards and
right-hand-side dependencies that may update an invariant variable. Reduction
claims require matching reference verdicts and a liftable unsafe witness; an
unsupported expression or disagreement fails the comparison. The retained
bundle records matching finite-fixture verdicts and measurable reduction, while
unavailable optional solver binaries contribute no equivalence evidence.

Authority-constrained planning uses authenticated operation catalogues,
immutable open-ended plan graphs, explicit bounded loops, typed continuations
and value bindings. The evaluation's dynamic-code mode emits an inert
\texttt{ModeledProgram} graph; it is not compiled, evaluated, or sent to an
executor. Every declared effect is re-authorised at modeled action time.
Security is a hard selection constraint;
eligible plans are then ordered by utility, authority footprint, sensitive
observations, irreversible effects, cost and a deterministic fingerprint.

\section{Results}
\textbf{Repository validation:} the retained Windows/CPython 3.12.12 run at
commit \texttt{6fe6b584500e} passed 220 tests with 90.25\% branch coverage,
13 schema checks, deterministic regeneration, static analysis, wheel build,
and installed-CLI smoke. The record is
\texttt{research/output/validation/6fe6b584500e/}.\\
\textbf{Semantic conformance:} the reference interpreter and, where the optional
binaries are present, the Z3 and nuXmv backends decide the same bounded IR
variants and agree on every decision, refuted variants included, and the
negative control still fails as designed. Table~\ref{tab:semantic-conformance}
counts them. A separate sweep closes the
population instead of listing it: the nine independent delegation defects have
512 subsets, and the Z3 and nuXmv encodings reproduce the reference interpreter's verdict
class on every one of them, with all three backends refuting 511. Agreement between
two encodings of one IR over a closed set of defect combinations is evidence
about the encoding, not a proof of the delegation semantics.\\
\textbf{Native SLED:} the retained bounded reproduction recovers all three
historical trace counts and their canonical-state compression; the canonical
fixtures are not directly comparable with the archived prototype, whose
enumeration differed. Table~\ref{tab:native-sled-reproduction} records both.\\
\textbf{SLED-MC/SLED-V:} the retained scaling bundle shows how far the COI
reduction shrinks the finite models and that no verdict is downgraded after
reduction; Table~\ref{tab:sledv-scaling} records the sizes and the checks. A
SLED-V reduction is bounded evidence on the tested finite models, not a proof
for unbounded deployments.\\
\textbf{Direction security models:} the retained checksummed bundle reports
safe exhaustion for the canonical finite disclosure and delegation models and
minimal witnesses for every seeded selector, disclosure, attribution, and
delegation defect. Delegation remains disabled at runtime.\\
\textbf{Cedar differential:} \emph{evaluation ready; Cedar unavailable and not
invoked}. The retained preflight contains translated PARC requests and
in-memory-oracle decisions but no parity result.\\
\textbf{Real models and AgentDojo:} the retained breadth evidence is counted in
Table~\ref{tab:agentdojo-breadth}, with the protocol identity of each bundle in
Table~\ref{tab:agentdojo-fingerprints}. Both tables are generated from the
retained result and protocol JSON. They record what the floor model did, not a
defence result: at this model scale utility is zero, so the security counts
carry no signal, and the standing limitation is stated in the table caption.
Real-model and AgentDojo claims stay bounded until a capable model completes
the protocols.

% GENERATED FILE -- do not edit by hand.
% Run: python scripts/generate_fourth_year_tables.py
% Source: research/output/runs/agentdojo-stratified-banking-v1/result.json SHA-256: 17edd75b1b5e6e03df835154457b60de3dff6c4f963020c7dd38c3bdd57777fa
% Source: research/output/runs/agentdojo-stratified-banking-v1/protocol.json SHA-256: cbf41d0ae0242311fadf67b2340a723cd06c20a26a3d11b1ec25ca5ade1e8cf1
% Source: research/output/runs/agentdojo-stratified-slack-v1/result.json SHA-256: 69e60b580d24bf9a28796d267e5434ec5afb9d664fe7dd697019fa47ccff2b15
% Source: research/output/runs/agentdojo-stratified-slack-v1/protocol.json SHA-256: 9c2ce6e181a17faaf96f061039b3cce5c9c209d965d17cd4951f7819a405946d
% Source: research/output/runs/agentdojo-stratified-travel-v1/result.json SHA-256: a4e85c02b7533a8ff71f536aca87b5a4594ad6467ff4828f5a99c004fc91a6d1
% Source: research/output/runs/agentdojo-stratified-travel-v1/protocol.json SHA-256: ee62fbc6ae0fbf0f4dcf6dcddc8e9e57a95db3693b67e8ec37c9498c5135fa5c
% Source: research/output/runs/agentdojo-stratified-workspace-v1/result.json SHA-256: 810e24057e8b961f49ae88c1b47901049b2443a5116c2057a0c099adcafb70ce
% Source: research/output/runs/agentdojo-stratified-workspace-v1/protocol.json SHA-256: 1d75d38896dc76d9d0a8071143f67fe0d05ade74acb620a4a8fd2cd80f6a711b
% Source: research/output/runs/agentdojo-pilot-banking-v1/result.json SHA-256: 7e74dfa894f89d13e1db9e35495b484c6e47eae4ad8bc4195608a79633983ab0
% Source: research/output/runs/agentdojo-pilot-banking-v1/protocol.json SHA-256: 44904393e0c13f582f1116bc1ed4c3c0227da55c87252ce50d5e7d709b0830a1
\begin{table}[t]
\centering
\small
\begin{tabular}{llrrrrl}
\toprule
\textbf{Run} & \textbf{Cells} & \textbf{Complete} & \textbf{Utility} & \textbf{Security} & \textbf{Incomplete} & \textbf{Failure categories} \\
\midrule
banking (\textit{stratified}, 4 tasks) & 24 & 24 & 0 & 24 & 0 & utility 24 \\
slack (\textit{stratified}, 5 tasks) & 30 & 24 & 0 & 24 & 6 & model 6, policy 6, utility 24 \\
travel (\textit{stratified}, 10 tasks) & 60 & 60 & 0 & 60 & 0 & policy 42, utility 60 \\
workspace (\textit{stratified}, 13 tasks) & 78 & 78 & 0 & 78 & 0 & policy 8, utility 78 \\
banking (\textit{pilot}, 2 tasks) & 12 & 12 & 0 & 12 & 0 & utility 12 \\
\textbf{Total} & \textbf{204} & \textbf{198} & \textbf{0} & \textbf{198} & \textbf{6} & \\
\bottomrule
\end{tabular}
\caption{Retained AgentDojo breadth evidence, counted from each bundle's own \texttt{result.json}. 204 cells, 198 complete, 0 with utility, 198 reporting no security violation. The model is a 0.5B floor configuration, so utility is zero and the security counts are vacuous: a model that completes no task cannot be shown to be protected. Failure counts are the categories the run recorded: each cell contributes one count per recorded failure, so one cell can carry several categories and a cell that completed can still carry a category such as \texttt{utility}. The six incomplete slack cells were later diagnosed as capability failures of that model, not adapter gaps; a module rerun of the same suite reports them as parser failures. Section 27 of the verification record holds the diagnosis and the repeatability measurement.}
\label{tab:agentdojo-breadth}
\end{table}


% GENERATED FILE -- do not edit by hand.
% Run: python scripts/generate_fourth_year_tables.py
% Source: research/output/runs/agentdojo-stratified-banking-v1/result.json SHA-256: 17edd75b1b5e6e03df835154457b60de3dff6c4f963020c7dd38c3bdd57777fa
% Source: research/output/runs/agentdojo-stratified-banking-v1/protocol.json SHA-256: cbf41d0ae0242311fadf67b2340a723cd06c20a26a3d11b1ec25ca5ade1e8cf1
% Source: research/output/runs/agentdojo-stratified-slack-v1/result.json SHA-256: 69e60b580d24bf9a28796d267e5434ec5afb9d664fe7dd697019fa47ccff2b15
% Source: research/output/runs/agentdojo-stratified-slack-v1/protocol.json SHA-256: 9c2ce6e181a17faaf96f061039b3cce5c9c209d965d17cd4951f7819a405946d
% Source: research/output/runs/agentdojo-stratified-travel-v1/result.json SHA-256: a4e85c02b7533a8ff71f536aca87b5a4594ad6467ff4828f5a99c004fc91a6d1
% Source: research/output/runs/agentdojo-stratified-travel-v1/protocol.json SHA-256: ee62fbc6ae0fbf0f4dcf6dcddc8e9e57a95db3693b67e8ec37c9498c5135fa5c
% Source: research/output/runs/agentdojo-stratified-workspace-v1/result.json SHA-256: 810e24057e8b961f49ae88c1b47901049b2443a5116c2057a0c099adcafb70ce
% Source: research/output/runs/agentdojo-stratified-workspace-v1/protocol.json SHA-256: 1d75d38896dc76d9d0a8071143f67fe0d05ade74acb620a4a8fd2cd80f6a711b
% Source: research/output/runs/agentdojo-pilot-banking-v1/result.json SHA-256: 7e74dfa894f89d13e1db9e35495b484c6e47eae4ad8bc4195608a79633983ab0
% Source: research/output/runs/agentdojo-pilot-banking-v1/protocol.json SHA-256: 44904393e0c13f582f1116bc1ed4c3c0227da55c87252ce50d5e7d709b0830a1
\begin{table}[t]
\centering
\small
\begin{tabular}{lllll}
\toprule
\textbf{Bundle} & \textbf{Protocol suite id} & \textbf{Source commit} & \textbf{Fingerprint} & \textbf{Model} \\
\midrule
agentdojo-stratified-banking-v1 & banking:expanded:user\_task\_0 & \texttt{1127281372dc} & \texttt{da3024092bc533e8} & Qwen/Qwen2.5-0.5B-Instruct \\
agentdojo-stratified-slack-v1 & slack:expanded:user\_task\_0 & \texttt{1127281372dc} & \texttt{936563e1782d1001} & Qwen/Qwen2.5-0.5B-Instruct \\
agentdojo-stratified-travel-v1 & travel:expanded:user\_task\_0 & \texttt{1127281372dc} & \texttt{54781fe7e475522b} & Qwen/Qwen2.5-0.5B-Instruct \\
agentdojo-stratified-workspace-v1 & workspace:expanded:user\_task\_0 & \texttt{1127281372dc} & \texttt{cbe8f9ad9a5f1b8b} & Qwen/Qwen2.5-0.5B-Instruct \\
agentdojo-pilot-banking-v1 & banking:expanded:user\_task\_0 & \texttt{1127281372dc} & \texttt{185539f9d2650759} & Qwen/Qwen2.5-0.5B-Instruct \\
\bottomrule
\end{tabular}
\caption{Protocol identity of the retained bundles. A repeat that changes the suite, task set, bounds, model identity, prompts, or inputs changes the fingerprint, so a matching fingerprint is what makes a rerun comparable; a matching fingerprint also implies the recorded source commit. The fingerprints are truncated to their first 16 hexadecimal digits here and kept whole in each bundle's \texttt{result.json}. The stratified bundles were retained before the sweep passed the selection's own suite id, so their suite ids use the earlier \texttt{<suite>:expanded:<first task>} spelling rather than the selection manifest's \texttt{<suite>:stratified-v1}; every id here is the one the run recorded.}
\label{tab:agentdojo-fingerprints}
\end{table}


% GENERATED FILE -- do not edit by hand.
% Run: python scripts/generate_fourth_year_tables.py
% Source: research/output/runs/three-way-agreement-v1/result.json SHA-256: bf3068da9078919a3a787aea911b4f4a85704138559b2dcbcc3553872db0639b
\begin{table}[t]
\centering
\small
\begin{tabular}{lr}
\toprule
\textbf{Quantity} & \textbf{Count} \\
\midrule
IR variants checked & 11 \\
Decided by Z3 & 11 \\
Decided by nuXmv & 11 \\
Refuted by the reference interpreter & 10 \\
Reference and Z3 agree & 11 \\
Reference and nuXmv agree & 11 \\
All three agree & 11 \\
Negative control & holds \\
\bottomrule
\end{tabular}
\caption{Semantic conformance over the same bounded IR variants: the reference interpreter, the optional Z3 backend and the optional nuXmv backend decide the same variants and agree on every decision, refuted variants included, and the negative control confirms the harness can still fail. Counts are read from the bundle's own \texttt{result.json}, so the table cannot drift from the evidence it cites.}
\label{tab:semantic-conformance}
\end{table}


% GENERATED FILE -- do not edit by hand.
% Run: python scripts/generate_fourth_year_tables.py
% Source: research/output/runs/native-sled-partb-reproduction-v1/result.json SHA-256: e3751abf62576d85333f08d8f771e56cb11141f195932d1daa5b79fd7d56158c
% Source: research/output/runs/native-sled-reproduction-v1/result.json SHA-256: de42dadd2480a66faf537abd744d117b817f5351c9240d5b368f434df0607657
\begin{table}[t]
\centering
\small
\begin{tabular}{lrrrl}
\toprule
\textbf{Environment} & \textbf{Expected traces} & \textbf{Reproduced} & \textbf{Match} & \textbf{Canonical states} \\
\midrule
env-01 & 422{,}535 & 422{,}535 & yes & 5 \\
env-02 & 996{,}451 & 996{,}451 & yes & 17 \\
env-03 & 43{,}621 & 43{,}621 & yes & 9 \\
\midrule
\textbf{Total} & 1{,}462{,}607 & 1{,}462{,}607 & yes & 31 \\
\midrule
\multicolumn{5}{l}{\emph{Canonical fixtures against the archived prototype}} \\
\multicolumn{5}{l}{Pairs compared: 3; negative controls killed: 5 of 5} \\
\multicolumn{5}{l}{Directly comparable with the archived trace claim: no (enumeration\_change)} \\
\multicolumn{5}{l}{Reproduced under depth bound 3; runtime is omitted from the deterministic fixture}} \\
\bottomrule
\end{tabular}
\caption{Native SLED reproduction. The historical environments reproduce their exact trace counts under the recorded depth bound, and the canonical-state compression collapses the raw traces to a handful of states. The archived prototype enumerated transitions differently, so its headline trace count is not directly comparable; this is reproduction and conformance, not an unbounded security proof. Counts are read from each bundle's own \texttt{result.json}.}
\label{tab:native-sled-reproduction}
\end{table}


% GENERATED FILE -- do not edit by hand.
% Run: python scripts/generate_fourth_year_tables.py
% Source: research/output/runs/sledv-scaling-v1/result.json SHA-256: 22f312a535a6ccb7163b1a6f86b6a0f928f1be1366baac5daab54d1d7461d81b
\begin{table}[t]
\centering
\small
\begin{tabular}{lrrrr}
\toprule
\textbf{Family} & \textbf{Fixtures} & \textbf{Largest model} & \textbf{After reduction} & \textbf{Reduction} \\
\midrule
noise scaling & 10 & 41{,}449 & 2 & 20724.5$\times$ \\
principal scaling & 8 & 2{,}517 & 2 & 1258.5$\times$ \\
depth scaling & 10 & 32 & 1 & 32.0$\times$ \\
\midrule
\multicolumn{5}{l}{Total fixtures: 28; COI-equivalent fixtures: 23} \\
\multicolumn{5}{l}{Verdict downgrades after reduction: none} \\
\multicolumn{5}{l}{Reference and Z3 agree: yes} \\
\bottomrule
\end{tabular}
\caption{SLED-V reduction over three scaling families. The reduction collapses the finite models by up to the factor shown, no verdict is downgraded after reduction, and the reference and Z3 decisions agree. Sizes and checks are read from the bundle's own \texttt{result.json}. This is bounded evidence on the tested finite models; it is not a proof of security for arbitrary or unbounded deployments.}
\label{tab:sledv-scaling}
\end{table}


No numerical claim is entered manually. Tables and figures are generated from
retained, versioned result JSON and include input checksums.

\section{Prioritised future work}

The immediate priority is operator-reviewed evidence: run the deliberately
small dual-backend laptop smoke, stop for review, then consider the full
planning and pinned AgentDojo protocols. Each live run must retain its exact
model and runtime identities, failures, native and Conflux metrics, and input
checksums. Paper tables must be generated solely from those result records.

The next policy milestone is executing the pinned Cedar differential corpus.
Successful offline request translation is readiness evidence only. Delegation
activation remains a later, separate gate and requires complete Cedar agreement
plus certificate, disclosure, attribution, expiry, revocation, and atomic-use
evidence.

The next security-model extensions are richer operation-specific argument
effects, any future activation of scoped attenuating delegation, and
origin-bound persistent-memory authority. They remain denied or conservative
until their semantics, revocation behavior, and counterexample properties are
specified. Formal work should then extend
the solver subset, relational properties, and independent implementation
conformance while keeping bounded and unbounded claims distinct.

\section{Limitations}
Conflux assumes sound authentication, provenance, policy state and complete
mediation. Authentication makes provenance accurate but does not grant
organisational authority; an authenticated external principal with no
relevant permissions will still cause ITES to block actions not jointly
authorised. It does not prevent actions already authorised for every influencing
Principal and does not guarantee subjective intent alignment. If both
influencing principals can perform an operation such as
\texttt{send\_email}, an attacker-controlled input may still influence which
recipient or attachment is selected---ITES confines authority, not harm
within authorised actions. Model-level defences and human judgement remain
relevant for deciding which already-authorised action is appropriate; ITES
does not guarantee semantic correctness of choices within the authority
envelope. Organisations may rely on human/model semantic judgement with
residual risk, use explicit delegation where intended, or disallow the
workflow. Conservative
contexts can over-approximate influence. Consider a workflow where Alice asks
the agent to read an email from an authenticated external supplier and create
a calendar event: the Principal Context becomes $\{\text{Alice},
\text{Supplier}\}$, and if Supplier lacks calendar permission, ITES blocks the
action. This is the intended zero-trust cost, not a bug. Delegation has a bounded model and
lifecycle representation but remains unsupported by operational ITES. Cedar
request translation has not been compared with a live Cedar binary.
Authorised-read safety is not noninterference. Formal results are scoped to
their model, backend and conformance evidence. Provider adapters are research
prototypes, not production isolation boundaries. Real-model, AgentDojo and
four-mode planning comparisons remain pending until gated runs produce retained
evidence; the raw upstream AgentDojo fixture validates translation only.

\section{Conclusion}
Conflux now has one tested security semantics, a runnable offline vertical
slice, argument-aware disclosure and attribution, authenticated bounded
planning, property-reduced verification, and a disabled scoped-delegation
foundation. The claim it defends is narrower than \emph{security for agents}:
an action is allowed only when every influencing principal's own permissions
allow it, and an input fetched from elsewhere cannot acquire authority it was
never granted, whatever the model proposes. The evidence in this paper is
bounded and reproducible. Eleven planned variants agree three ways under two
independent backends, the archived prototype reproduces its published trace
count under a stated enumeration change, and property reduction keeps the
scaling advantage measured on the fixtures exercised here. The controls hold and
the ten variants that are false are reported as refuted, which is what the
method exists to produce: refutation is a result, not a failure of the run. The
remaining research contribution is stronger evidence and carefully gated
activation: real-model, external-benchmark, Cedar-parity and formal artefacts
whose claims remain tied to explicit assumptions, models, and bounds. The
limitations section states where that evidence does not yet reach.

\bibliographystyle{plainnat}
\bibliography{references}
\end{document}
